\scanpage{320}
\mspage{236}
% archive sheet number: 319

2\textsuperscript{o}) Notons aussi que dans la condition b'), on ne peut supprimer l'hypothèse que $I'$ satisfait $\mathrm{PF}2$ (en plus de $I'$ $0$-connexe et les $ {}_{i'}\backslash I$ filtrantes). Pour le voir, prenons \deleted{$f$} $I'=I$, $f=\mathrm{id}$, donc l'hyp. \unsure{b')|a)} $(\Leftrightarrow b\Leftrightarrow b')$ signifie simplement que $I$ est filtrante, $f$ étant bien sûr cofinal.

Mais b') affaiblie signifie que $I$ est $0$-connexe et les $ {}_{i}\backslash I$ sont filtrantes, et on sait que ceci n'implique \uline{pas} que $I$ soit filtrante, $=$ si on suppose $I$ pseudo-filtrante. Un \deleted{exemple} simple est donné par $I$ définie par

\[
\mathrm{Ob}\,I=\mathbb{N},\qquad
\deleted{\unsure{\mathrm{Fl}}}\deleted{\unsure{\mathrm{Ens}}}
\mathrm{Hom}_{I}(i,j)=
\begin{cases}
\varnothing&\text{si }i>j\\
\{\mathrm{id}_i\}&\text{si }i=j\\
\{\pm1\}_{\unsure{i,j|j,i}}&\text{si }i<j,
\end{cases}
\]

la composition des flèches se faisant par la multiplication des signes $\unsure{\varepsilon|\epsilon}_{k,j}\varepsilon'_{ji}=(\varepsilon\unsure{\bar{\varepsilon}'|\varepsilon'})_{ki}$. (en posant $\mathrm{id}_i=(+1)_{i,i}$). Les catégories $ {}_i\backslash I$ sont \inserted{\unsure{isom.} entre elles et sont} \uline{ordonnées} et ps filtr. donc filtrantes, \deleted{\unsure{d'ailleurs}} puis $I$ est ps. filtrante, \unsure{pour}\deleted{\unsure{tant}} elle n'est pas filtrante, car ne satisfait pas $\mathrm{PF}\unsure{2|?}$.

3) D'autre part, si $f$ n'est pas fidèle, alors c) (équivalente à \uline{b'')}) n'implique pas a), comme on voit en prenant $I'=e$ ; alors b'') équivaut à $I$ ps. filtrante, qui n'implique pas que $I$ soit filtrante.

